\section{Detection of zeros from interval bounds}\label{sec:window}

The estimate in Proposition~\ref{prop:arithmetic} is one-sided
and is available only on selected intervals. Why can these intervals exclude a zero even when nothing comparable is known in the
gaps? For a finite family with positive weights the answer comes from choosing the
order of differentiation after the interval is known. Taking many derivatives of a Laplace
transform concentrates its integral on one of the intervals where the bound holds.
Comparing a real derivative with one at a complex point makes
the kernel in the upper bound nonnegative. The deduction in
Section~\ref{sec:deduction} applies the criterion with
intervals \([Y,Y^{3/2}]\) and an exponent chosen below the real part of the
proposed zero.

\begin{definition}[The transforms of a weighted family]\label{def:window-transforms}
Let \((F,w)\) be a finite positively weighted family as in
Definition~\ref{def:family-weights}. In logarithmic coordinates set
\(g_{F,w}(u)=G_{F,w}(e^u)\) for \(u\ge0\).
For a primitive character \(\chi\), let \(H_\chi\) be the Mellin
transform of its centered sum. Let \(\Phi_{F,w}\) be the Laplace
transform of the real-valued function \(g_{F,w}\). Explicitly,
\[
 H_\chi(s)=\int_1^\infty G_\chi(y)y^{-s-1}\,dy,\qquad
 \Phi_{F,w}(s)=\int_0^\infty g_{F,w}(u)e^{-su}\,du
                       \qquad(\RePart s>1).
\]
Absolute convergence follows from the growth estimate at the start
of the proof below. With the family fixed, write \(g=g_{F,w}\),
\(\Phi=\Phi_{F,w}\), and \(w_\chi=w(\chi)\).
\end{definition}

\begin{theorem}[A zero-free half-plane from bounds on intervals]\label{thm:window}
Let \((F,w)\) be as in Definition~\ref{def:window-transforms},
with \(G_\chi\) as in \eqref{eq:centered-sum}.
Suppose that \(0<\theta<1\), \(\delta>0\), \(C>0\), and
\[
 \forall Z\ge1\quad\exists Y>Z\quad\forall y\in[Y,Y^{1+\delta}],
             \qquad G_{F,w}(y)\le Cy^\theta .
\]
Then each \(\chi\in F\) satisfies
\[
 \RePart s>\theta,\quad (\chi,s)\ne(1,1)
         \quad\Longrightarrow\quad L_{f_\chi}(s,\chi)\ne0.
\]
\end{theorem}

\begin{proof}
\proofstep{The poles of the transform.}
For every \(y\ge1\), absolute values in the defining sum give
\[
 |S_\chi(y)|\le y\sum_{n\le y}\frac{\log n}{n}
             \ll y(1+\log^2 y).
\]
The centering terms satisfy the same bound for each fixed character.
Thus the real function \(g(u)=G_{F,w}(e^u)\) satisfies \(|g(u)|\le B e^{2u}\) on
\(u\ge0\), for a constant \(B\) depending on the weighted family.

For \(\RePart s>1\) the Mellin transform of \(S_\chi\) is
\[
 \int_1^\infty S_\chi(y)y^{-s-1}\,dy
    =-\frac{L'(s,\chi)}{s(s-1)L(s,\chi)} .
\]
To justify the exchange with the Dirichlet series, first integrate
a single summand over \(y\in[n,\infty)\).
\[
 \int_n^\infty(y/n-1)y^{-s-1}\,dy
   = n^{-s}\left(\frac1{s-1}-\frac1s\right).
\]

The same calculation with \(s\) replaced by \(\RePart s>1\)
majorizes the absolute integrals, since
\(\sum_n\Lambda(n)n^{-\RePart s}<\infty\).
Transforming the two centering terms of \eqref{eq:centered-sum} therefore gives
\begin{align*}
 H_\chi(s)&=-\frac{L'(s,\chi)}{s(s-1)L(s,\chi)}
        +\frac{L'(1,\chi)}{(s-1)L(1,\chi)} \qquad(\chi\ne 1),\\
 H_1(s)&=-\frac{\zeta'(s)}{s(s-1)\zeta(s)}
        -\frac{1}{(s-1)^2}+\frac{1+\gamma}{s-1}.
\end{align*}
These formulas continue meromorphically to \(\RePart s>0\).
By Lemma~\ref{lem:boundary}, for \(\chi\ne1\) the coefficient of
\((s-1)^{-1}\) in the first term of \(H_\chi\) is
\(-L'(1,\chi)/L(1,\chi)\), and the second term cancels it.

For the conductor-one character one writes
\[
 \zeta(s)-\frac1{s-1}
   =\sum_{n\ge1}\left(n^{-s}-\int_n^{n+1}v^{-s}\,dv\right).
\]
On a compact subset of \(\RePart s>0\) each summand is bounded
by \(\sup_{s}|s|\,n^{-1-\inf_s\RePart s}\) by integration of
the derivative of \(v^{-s}\). At \(s=1\) the partial sums are
\(\sum_{n\le M}1/n-\log(M+1)\to\gamma\). Consequently, with
\(h=s-1\),
\[
 -\frac{\zeta'(1+h)}{(1+h)h\,\zeta(1+h)}
      =\frac{1}{h^2}-\frac{1+\gamma}{h}+O(1).
\]
These two terms cancel the two explicit centering terms in \(H_1\).

Substituting \(y=e^u\) in Definition~\ref{def:window-transforms} gives the
following continuation of \(\Phi\) from \(\RePart s>1\) to \(\RePart s>0\):
\[
 \Phi(s)=\tfrac12\sum_{\chi\in F}w_\chi
                         (H_\chi(s)+H_{\bar\chi}(s)).
\]
Here \(G_{\bar\chi}(y)=\overline{G_\chi(y)}\) on the real axis.
For \(0<\RePart\rho<1\) let \(m_\chi(\rho)\) be the order of
vanishing of \(L(s,\chi)\) at \(\rho\), with value zero if
\(L(\rho,\chi)\ne0\). Factoring out \((s-\rho)^{m_\chi(\rho)}\)
gives
\[
 \operatorname{Res}_{s=\rho}\Phi(s)=
 -\frac{1}{2\rho(\rho-1)}
       \sum_{\chi\in F}w_\chi
                \bigl(m_\chi(\rho)+m_{\bar\chi}(\rho)\bigr).
\]
If any order on the right is positive, the positive weights
make this a nonzero residue. The residue is real and positive
at a real zero, since \(\rho(1-\rho)>0\) there. The Euler products and Lemma~\ref{lem:boundary}
exclude all other poles in \(\RePart s\ge 1\), and the centering
has removed the singularity at \(1\).

\proofstep{Concentrating the derivative integrals.}
The residue calculation identifies what a zero contributes. To test for that
contribution using bounds on only one interval at a time, fix a real \(\sigma>2\), to
be chosen below. For a nonnegative
integer \(k\) and a complex point \(z\) with \(\RePart z=\sigma\),
put \(T_k(z)=(-1)^k\Phi^{(k)}(z)/k!\).
With the growth constant \(B\) fixed at the start of the proof, differentiation
under the integral is justified by
\[
  |u^kg(u)e^{-zu}|\le B u^k e^{-(\sigma-2)u},\qquad
  \int_0^\infty u^k e^{-(\sigma-2)u}\,du
       =\frac{k!}{(\sigma-2)^{k+1}} .
\] Write \(g^+=\max(g,0)\) for the positive part.
For a real \(\tau\), choose \(c\in\C\) with \(|c|=1\) and
\(\RePart(cT_k(\sigma+i\tau))=|T_k(\sigma+i\tau)|\).
Since \(g\) is real, the normalized derivatives at \(\sigma\) and \(\sigma+i\tau\) satisfy
\[
 |T_k(\sigma+i\tau)|+\RePart T_k(\sigma)
  =\int_0^\infty g(u)\frac{u^k e^{-\sigma u}}{k!}
                        (1+\RePart(ce^{-i\tau u}))\,du .
\]
This kernel lies between zero and twice the real exponential kernel,
which gives
\begin{equation}\label{eq:two-centers}
 |T_k(\sigma+i\tau)|+\RePart T_k(\sigma)
 \leq 2\int_0^\infty g^+(u)\frac{u^k}{k!}e^{-\sigma u}\,du
 \qquad (\tau\in\R).
\end{equation}

Writing \(U=\log Y\), set \(W=[U,(1+\delta)U]\).
On \(W\), \(g^+(u)\leq C e^{\theta u}\). The integral over \(W\)
is at most \(C\delta U/(\sigma-\theta)^k\) because
\(v^k e^{-v}/k!\leq 1\) for \(v\geq 0\).

Set \(a=\sigma-2\), \(b=1+\delta/2\), and \(k=\lceil ab U\rceil\).
Let \(X\) have the gamma distribution with density \(v^ke^{-v}/k!\)
for \(v>0\). On the complement of \(W\), the substitution \(v=au\)
and the bound \(|g(u)|\le B e^{2u}\) give the upper bound
\[
 B a^{-k-1}\Pr\{X\notin[aU,(1+\delta)aU]\},
 \qquad X\sim\operatorname{Gamma}(k+1,1).
\]

Fix \(q_-,q_+\) with \(1/b<q_-<1<q_+<(1+\delta)/b\).
As \(U\to\infty\), \(aU/(k+1)\to1/b\) and
\((1+\delta)aU/(k+1)\to(1+\delta)/b\). Hence the event
in the probability is contained in
\[
 \{X\le q_-(k+1)\}\cup\{X\ge q_+(k+1)\}
\]
for all large \(U\). 
In the first event take \(t=1-1/q_-<0\),
and in the second take \(t=1-1/q_+>0\). In both cases
Markov's inequality applied to \(e^{tX}\), whose mean is
\((1-t)^{-k-1}\), gives the bound
\[
 e^{-tq(k+1)}(1-t)^{-k-1}
          =e^{-(k+1)(q-1-\log q)} .
\]

Put \(c_\delta=\min_{q\in\{q_-,q_+\}}(q-1-\log q)>0\), and
choose \(a\) so that \(ab>1\) and \(\log((a+2)/a)<c_\delta\).
Thus \(a^{-k-1}e^{-c_\delta k}\le a^{-1}\sigma^{-k}\).
With \(D=\sigma-\theta\), \eqref{eq:two-centers} gives
\begin{equation}\label{eq:derivative-upper}
 |T_k(\sigma+i\tau)|+\RePart T_k(\sigma)
       \ll (1+U)D^{-k}.
\end{equation}
The constant is uniform in \(\tau,k,U\) along the chosen sequence of intervals.

\begin{construction}
\begin{tabular}{@{}l@{\hspace{1.2em}}l@{}}
\toprule
\textit{Part of the integral} & \textit{Bound used} \\
\midrule
\(u\in[U,(1+\delta)U]\) & \(g^+(u)\le Ce^{\theta u}\) \\
\(u\notin[U,(1+\delta)U]\) & \(|g(u)|\le Be^{2u}\), gamma tail in \(au\) \\
Both centers \(\sigma,\sigma+i\tau\) & \(0\le 1+\RePart(ce^{-i\tau u})\le2\) \\
\bottomrule
\end{tabular}
\caption{Where the derivative estimate comes from in the proof of Theorem~\ref{thm:window}. Here \(U=\log Y\), \(a=\sigma-2\) and \(k=\lceil a(1+\delta/2)U\rceil\).}\label{fig:analytic-kernel}
\end{construction}

\proofstep{Derivatives forced by a pole.}
We now have an upper bound at the derivative orders supplied by each interval. To
obtain a lower bound at those same orders, suppose that \(\Phi\) has a pole with real
part greater than \(\theta\).
First consider the case in which the pole \(\beta\) nearest to \(\sigma\)
is unique, real and at distance \(d=\sigma-\beta<D\).
Its residue \(r\) is positive. Choose \(d<R<D\) so that the closed disk
\(|s-\sigma|\le R\) contains no other pole. Subtract the principal part
and set \(H(s)=\Phi(s)-r/(s-\beta)\). With \(M=\max_{|s-\sigma|=R}|H(s)|\),
Cauchy's estimate gives \(|H^{(k)}(\sigma)|/k!\le M/R^k\). Consequently
\[
       2\RePart T_k(\sigma)\ge 2r d^{-k-1}-2M/R^k.
\]

At \(\tau=0\), \(T_k(\sigma)\) is real since \(g\) is real.
The left side in \eqref{eq:derivative-upper} dominates
\(2\RePart T_k(\sigma)\). Multiplying both estimates by \(d^k\)
gives a positive limit inferior on the lower side and an upper side
bounded by a constant times \((1+U)(d/D)^{\lceil abU\rceil}\to0\).

For the remaining cases write
\(d_\sigma=\inf_\rho|\sigma-\rho|\), the distance from the
real center to the pole set. We seek a \(z=\sigma+i\tau\)
with a unique nearest pole at distance \(d<\min(D,d_\sigma)\).

If the open disk of
radius \(D\) at \(\sigma\) contains no pole, align \(\tau\) with
the imaginary part of any pole whose real part is above \(\theta\).
Otherwise a nearest pole \(\rho\) is nonreal. Setting
\(\tau=\operatorname{Im}\rho\) gives
\[
        |z-\rho|=\sigma-\RePart\rho
        <\sqrt{(\sigma-\RePart\rho)^2+(\operatorname{Im}\rho)^2}
        =d_\sigma<D .
\]

In either case the inequalities are strict and persist
under small perturbations of \(\tau\). On a small closed vertical segment of possible centers, all such
disks lie in a compact subset of \(\RePart s\ge\theta>0\).
They can meet only finitely many poles. For distinct poles
\(a_1+ib_1,a_2+ib_2\), equality of squared distances is the equation
\[
    (\sigma-a_1)^2-(\sigma-a_2)^2
       +b_1^2-b_2^2-2\tau(b_1-b_2)=0 .
\]
If \(b_1=b_2\), the constant is nonzero because both \(a_i<\sigma\)
and \(a_1\ne a_2\). Thus finitely many values of \(\tau\) exhaust
all ties. Choosing \(\tau\) off this set preserves the strict
distance bounds and gives a unique nearest pole.

Write \(\rho\) for the closest pole to \(z\), with residue \(r\ne 0\).
Choose a radius \(R\in(d,D)\) below the distance from \(z\) to any other
pole and below the distance from \(\sigma\) to every pole.
Both closed disks of radius \(R\) lie in \(\RePart s>\theta\).
Cauchy's estimate for the regular part at \(z\) and for \(\Phi\)
at \(\sigma\) gives, for a constant \(M\) independent of \(k\),
\[
 |T_k(z)|+\RePart T_k(\sigma)\geq |r|d^{-k-1}-M/R^k .
\]
Multiplying the lower and upper bounds by \(d^k\) gives
\[
 \frac{|r|}{d}-M(d/R)^k
         \ll (1+U)(d/D)^k .
\]
The left side tends to \(|r|/d>0\). Since \(d/D<1\) and
\(k\ge abU\), the right side tends to zero.
This excludes the pole and hence the zero.
\end{proof}
